\part{Ketten}

\chapter{Die Irrfahrt}

\section{Ein Prozess ist ein Pfad}

\begin{definition}
\boxtitelm{Definition 1.1 --- stochastischer Prozess}
Ein \Begriff{stochastischer Prozess} ist eine Familie von Zufallsgrößen
\[ \bigl(X_t\bigr)_{t \in T} \]
auf demselben Wahrscheinlichkeitsraum. Für festes $\omega$ heißt die Abbildung $t \mapsto X_t(\omega)$ ein \Begriff{Pfad}.

Ist $T = \{0,1,2,\ldots\}$, spricht man von \Begriff{diskreter}, bei $T = [0,\infty)$ von \Begriff{stetiger Zeit}.
\end{definition}

\begin{merksatz}
\boxtitel{Zwei Blickrichtungen}
\begin{center}
{\small
\renewcommand{\arraystretch}{1.52}
\begin{tabular}{@{}P{2.4cm}P{6.2cm}@{}}
\toprule
\textbf{Verteilungen} & Wie ist $X_t$ für festes $t$ verteilt? Was ist die gemeinsame Verteilung von $X_{t_1},\ldots,X_{t_n}$?\\
\textbf{Pfade} & Ist der Verlauf stetig? Monoton? Wie oft überquert er die Null? Wann erreicht er erstmals einen Wert?\\
\bottomrule
\end{tabular}}
\end{center}

\begin{center}
\begin{minipage}{0.9\textwidth}
\textbf{Die Verteilungen sind leicht, die Pfade sind schwer --- und fast alle interessanten Fragen betreffen die Pfade.}
\end{minipage}
\end{center}

\begin{center}
\begin{minipage}{0.88\textwidth}
\kursiv{Die Verteilung von $X_t$ zu kennen, sagt nichts über die Verweildauer oberhalb der Null oder über die Zahl der Rückkehrpunkte. \textbf{Genau diese Größen entscheiden aber darüber, wie ein Verlauf auf den Betrachter wirkt} --- ob er wie Zufall aussieht oder wie ein Trend. Dieses Kapitel zeigt es an der einfachsten Kette überhaupt.}
\end{minipage}
\end{center}
\end{merksatz}

\begin{definition}
\boxtitelm{Definition 1.2 --- die einfache Irrfahrt}
\[ S_0 = 0, \qquad S_n = \sum_{k=1}^{n} \xi_k \]
mit unabhängigen $\xi_k$, die die Werte $+1$ und $-1$ mit Wahrscheinlichkeit je $\tfrac12$ annehmen.
\end{definition}

\begin{merksatz}
\boxtitel{Was leicht ist}
\begin{center}
{\small
\renewcommand{\arraystretch}{1.5}
\begin{tabular}{@{}P{3.0cm}P{3.0cm}@{}}
\toprule
$\mathbb{E}[S_n]$ & $0$\\
$\operatorname{Var}[S_n]$ & $n$\\
$\mathbb{E}\bigl[|S_n|\bigr]$ & $\approx \sqrt{2n/\pi}$\\
$\Pr[S_{2n} = 0]$ & $\binom{2n}{n}2^{-2n} \approx \frac{1}{\sqrt{\pi n}}$\\
\bottomrule
\end{tabular}}
\end{center}

\kursiv{Bei $n = 2000$:} $\sqrt{2\cdot2000/\pi} = 35{,}7$; die Simulation ergab $35{,}6$.

\begin{center}
\begin{minipage}{0.88\textwidth}
\kursiv{\textbf{Die Wurzel ist die Signatur des Zufalls}: Nach zehntausend Schritten ist man typischerweise achtzig Schritte vom Ursprung entfernt, nicht zehntausend und nicht null. Alle vier Zeilen folgen aus Band 4 und brauchen keine Prozesstheorie. Der Rest dieses Kapitels braucht sie.}
\end{minipage}
\end{center}
\end{merksatz}

\section{Das Gesetz vom Arkussinus}

\begin{merksatz}
\boxtitel{Die falsche Erwartung}
Zwei Spieler werfen zweitausendmal eine Münze; nach jedem Wurf zahlt der Verlierer einen Euro. Wie lange liegt der eine vorn?

\begin{center}
\begin{minipage}{0.9\textwidth}
\textbf{Die meisten erwarten: ungefähr die Hälfte der Zeit jeder. Das ist der \kursiv{unwahrscheinlichste} aller Verläufe.}
\end{minipage}
\end{center}

\schritt{Satz 1.3 --- Lévy, 1939}
Der Anteil der Zeit, den die Irrfahrt oberhalb der Null verbringt, hat für großes $n$ die Verteilungsfunktion
\[ F(x) = \frac{2}{\pi}\arcsin\sqrt{x}, \qquad 0 \le x \le 1 . \]

\begin{center}
\begin{tikzpicture}[scale=1.0]
\draw[grau!60, line width=0.6pt] (0,0) -- (6.2,0);
\draw[grau!60, line width=0.6pt] (0,0) -- (0,2.4);
\draw[tief, line width=1.4pt, smooth, samples=80, domain=0.012:0.988] plot ({6*\x},{0.52/(3.1416*sqrt(\x*(1-\x)))*2.4});
\node[font=\scriptsize, color=grau, anchor=north] at (0,-0.1) {$0$};
\node[font=\scriptsize, color=grau, anchor=north] at (3,-0.1) {$\tfrac12$};
\node[font=\scriptsize, color=grau, anchor=north] at (6,-0.1) {$1$};
\node[font=\scriptsize\itshape, color=grau, anchor=west] at (2.0,2.1) {Dichte $\dfrac{1}{\pi\sqrt{x(1-x)}}$};
\end{tikzpicture}
\end{center}

\begin{center}
\begin{minipage}{0.9\textwidth}
\textbf{Die Dichte hat ihr Minimum bei $\tfrac12$ und wächst zu beiden Rändern hin über alle Grenzen.}
\end{minipage}
\end{center}
\end{merksatz}

\begin{merksatz}
\boxtitel{An zwanzigtausend Läufen geprüft}
Je $2000$ Schritte:

\begin{center}
{\small
\renewcommand{\arraystretch}{1.46}
\begin{tabular}{@{}P{2.6cm}P{2.2cm}P{2.4cm}@{}}
\toprule
\textbf{Zeitanteil oben} & \textbf{simuliert} & \textbf{theoretisch}\\
\midrule
unter $10\,\%$ & $0{,}214$ & $0{,}205$\\
$10$--$20\,\%$ & $0{,}090$ & $0{,}090$\\
$40$--$50\,\%$ & $0{,}062$ & $0{,}064$\\
\midrule
unter $10$ \kursiv{oder} über $90\,\%$ & $\mathbf{0{,}408}$ & $\mathbf{0{,}410}$\\
zwischen $40$ und $60\,\%$ & $\mathbf{0{,}128}$ & $\mathbf{0{,}128}$\\
\bottomrule
\end{tabular}}
\end{center}

\begin{center}
\begin{minipage}{0.9\textwidth}
\textbf{In zwei von fünf Spielen liegt einer der beiden über neunzig Prozent der Zeit vorn. Der ausgeglichene Verlauf tritt nur in jedem achten Spiel auf.}
\end{minipage}
\end{center}

\begin{center}
\begin{minipage}{0.88\textwidth}
\kursiv{\textbf{Und der Führende ist nicht der bessere Spieler} --- das Spiel ist fair. Wer solche Verläufe als Zeitreihe sieht, wird unweigerlich Phasen erkennen, Trends benennen und Erklärungen suchen. Es gibt keine. Dasselbe gilt für Börsenkurse, Krankheitsverläufe und Torfolgen: \textbf{Lange einseitige Strecken sind kein Hinweis auf eine Ursache, sondern die Regel.}}
\end{minipage}
\end{center}

\schritt{Und die Vorzeichenwechsel}
\begin{center}
{\small
\renewcommand{\arraystretch}{1.46}
\begin{tabular}{@{}P{3.2cm}P{2.0cm}P{2.2cm}@{}}
\toprule
& \textbf{simuliert} & $\sqrt{2n/\pi}$\\
\midrule
Nullstellen in $2000$ Schritten & $34{,}8$ & $35{,}7$\\
\bottomrule
\end{tabular}}
\end{center}

\kursiv{Nicht tausend, wie man bei einer fairen Münze vermuten könnte, sondern fünfunddreißig --- und die verteilen sich ungleichmäßig über den Verlauf.}
\end{merksatz}

\section{Rückkehr}

\begin{merksatz}
\boxtitel{Satz 1.4 --- Pólya, 1921}
Die einfache Irrfahrt kehrt in einer und in zwei Dimensionen mit Wahrscheinlichkeit $1$ zum Ausgangspunkt zurück, in drei Dimensionen nur mit Wahrscheinlichkeit
\[ 1 - \frac{1}{u} \approx 0{,}3405 . \]

\schritt{Der Gedanke des Beweises}
\kursiv{Sei $u$ die erwartete Zahl der Besuche im Ursprung} (den Start mitgezählt), und $q$ die Wahrscheinlichkeit, überhaupt zurückzukehren. Dann ist die Zahl der Besuche geometrisch verteilt:
\[ u = \frac{1}{1-q} \quad\Longleftrightarrow\quad q = 1 - \frac1u . \]

\kursiv{Andererseits} ist $u = \sum_{n\ge0} \Pr[S_n = 0]$, und für große $n$ gilt
\[ \Pr[S_{2n} = 0] \ \sim\ C\,n^{-d/2} . \]

\kursiv{Die Reihe} $\sum n^{-d/2}$ \kursiv{divergiert für $d \le 2$ und konvergiert für $d \ge 3$} (Band 3, Kapitel~9).

\begin{center}
{\small
\renewcommand{\arraystretch}{1.5}
\begin{tabular}{@{}P{1.6cm}P{2.4cm}P{2.6cm}@{}}
\toprule
$d$ & $u$ & $q$\\
\midrule
$1$, $2$ & $\infty$ & $1$\\
$\ge 3$ & endlich & $< 1$\\
\bottomrule
\end{tabular}}
\end{center}
$\blacksquare$ \kursiv{Die Berechnung von $u = 1{,}5164$ für $d = 3$ durch ein dreifaches Integral übergehen wir --- die erste Auslassung dieses Bandes.}

\schritt{Simuliert}
\begin{center}
{\small
\renewcommand{\arraystretch}{1.46}
\begin{tabular}{@{}P{1.6cm}P{3.0cm}P{2.2cm}@{}}
\toprule
$d$ & \textbf{Rückkehr in $5000$ Schritten} & \textbf{Grenzwert}\\
\midrule
$1$ & $0{,}990$ & $1$\\
$2$ & $0{,}726$ & $1$\\
$3$ & $\mathbf{0{,}332}$ & $\mathbf{0{,}3405}$\\
\bottomrule
\end{tabular}}
\end{center}

\begin{center}
\begin{minipage}{0.88\textwidth}
\kursiv{In der Ebene nähert sich der Anteil der Eins nur sehr langsam --- die erwartete Rückkehrzeit ist dort unendlich, obwohl die Rückkehr sicher ist. \textbf{Man kehrt gewiss zurück, aber man kann nicht sagen, wann}; das ist der Unterschied zwischen \enquote{fast sicher} und \enquote{im Mittel bald}, und er zieht sich durch den ganzen Band.}
\end{minipage}
\end{center}

\begin{center}
\begin{minipage}{0.88\textwidth}
\kursiv{Der Satz hat eine Deutung, die \Person{Shizuo Kakutani} zugeschrieben wird: \textbf{Ein betrunkener Mann findet nach Hause, ein betrunkener Vogel nicht.}}
\end{minipage}
\end{center}
\end{merksatz}

\section{Der Bruchfall}

\begin{warnung}
\boxtitelr{Was an der Irrfahrt besonders ist}
\begin{center}
{\small
\renewcommand{\arraystretch}{1.52}
\begin{tabular}{@{}P{2.6cm}P{6.0cm}@{}}
\toprule
\textbf{Symmetrie} & Schon eine minimale Drift ändert alles: Bei $p = 0{,}51$ ist die Rückkehr nicht mehr sicher.\\
\textbf{Unabhängigkeit} & Reale Verläufe haben Gedächtnis --- Kurse, Krankheiten, Populationen.\\
\textbf{Gleiche Schritte} & Bei Schritten mit unendlicher Varianz gelten weder die Wurzel noch der Arkussinus.\\
\bottomrule
\end{tabular}}
\end{center}

\schritt{Zur ersten Zeile}
Bei Aufwärtswahrscheinlichkeit $p > \tfrac12$ beträgt die Rückkehrwahrscheinlichkeit
\[ q = \frac{1-p}{p}\cdot\ldots = \min\left\{1,\ \frac{1-p}{p}\right\}\cdot 2p = 2(1-p) \ \text{ für } p \ge \tfrac12 . \]
\begin{center}
{\small
\renewcommand{\arraystretch}{1.46}
\begin{tabular}{@{}P{1.8cm}P{2.4cm}@{}}
\toprule
$p$ & $q$\\
\midrule
$0{,}500$ & $1{,}000$\\
$0{,}510$ & $0{,}980$\\
$0{,}550$ & $0{,}900$\\
\bottomrule
\end{tabular}}
\end{center}

\begin{center}
\begin{minipage}{0.88\textwidth}
\kursiv{\textbf{Ein Prozent Vorteil je Schritt genügt, um die sichere Rückkehr zu zerstören} --- und in zwanzig von tausend Fällen entfernt sich der Verlauf dann für immer. Die Symmetrie der einfachen Irrfahrt ist also kein harmloser Sonderfall, sondern die Grenze zwischen zwei Welten. Kapitel~8 zeigt, dass genau an dieser Grenze die Martingale stehen.}
\end{minipage}
\end{center}
\end{warnung}

\section{Übungsaufgaben}

\begin{aufgabe}
\boxtitel{Aufgabe 1.1 --- Die Wurzel nachrechnen}
Berechnen Sie $\Pr[S_{2n} = 0]$ für $n = 1, 2, 10$ und vergleichen Sie mit der Näherung $1/\sqrt{\pi n}$.

\loesung
\schritt{Schritt 1 --- Die exakte Formel}
\[ \Pr[S_{2n}=0] = \binom{2n}{n}\,2^{-2n} \]

\schritt{Schritt 2 --- Die Werte}
\begin{center}
{\small
\renewcommand{\arraystretch}{1.48}
\begin{tabular}{@{}P{1.2cm}P{2.6cm}P{1.8cm}P{2.0cm}@{}}
\toprule
$n$ & \textbf{exakt} & \textbf{Wert} & $1/\sqrt{\pi n}$\\
\midrule
$1$ & $2/4$ & $0{,}5000$ & $0{,}5642$\\
$2$ & $6/16$ & $0{,}3750$ & $0{,}3989$\\
$10$ & $184\,756/2^{20}$ & $0{,}1762$ & $0{,}1784$\\
\bottomrule
\end{tabular}}
\end{center}

\schritt{Schritt 3 --- Die Herkunft der Näherung}
Mit der Stirlingschen Formel aus Band 15, Kapitel~4:
\[ \binom{2n}{n} \approx \frac{4^{n}}{\sqrt{\pi n}} \quad\Longrightarrow\quad \Pr[S_{2n}=0] \approx \frac{1}{\sqrt{\pi n}} \]

\schritt{Schritt 4 --- Die Folge}
\[ \sum_{n\ge1} \frac{1}{\sqrt{\pi n}} = \infty \]

\begin{center}
\begin{minipage}{0.88\textwidth}
\kursiv{\textbf{Die Divergenz dieser Reihe ist der ganze Beweis von Satz 1.4 in einer Dimension.} Die erwartete Zahl der Nullbesuche ist unendlich, also ist die Rückkehrwahrscheinlichkeit eins. Dass eine so harmlose Reihe über die Wiederkehr eines Verlaufs entscheidet, ist einer der schönsten Übergänge der Wahrscheinlichkeitsrechnung.}
\end{minipage}
\end{center}
\end{aufgabe}

\begin{aufgabe}
\boxtitel{Aufgabe 1.2 --- Das Spiegelungsprinzip}
Zeigen Sie: Die Wahrscheinlichkeit, dass die Irrfahrt in $n$ Schritten das Niveau $a > 0$ erreicht, beträgt
\[ \Pr\bigl[\max_{k\le n} S_k \ge a\bigr] = 2\,\Pr[S_n > a] + \Pr[S_n = a] . \]

\loesung
\schritt{Schritt 1 --- Der Gedanke}
Man betrachte einen Pfad, der $a$ erreicht und bei $S_n < a$ endet. \kursiv{Man spiegele alles nach dem ersten Erreichen von $a$ an der Höhe $a$.}

\schritt{Schritt 2 --- Die Zuordnung}
\begin{center}
{\small
\renewcommand{\arraystretch}{1.5}
\begin{tabular}{@{}P{3.4cm}P{4.6cm}@{}}
\toprule
\textbf{Ursprünglicher Pfad} & erreicht $a$, endet bei $S_n = a - c$\\
\textbf{Gespiegelter Pfad} & erreicht $a$, endet bei $S_n = a + c$\\
\bottomrule
\end{tabular}}
\end{center}
Die Spiegelung ist eine Eins-zu-eins-Zuordnung, und beide Pfade haben dieselbe Wahrscheinlichkeit --- $2^{-n}$, denn jeder Pfad gleicher Länge ist gleich wahrscheinlich.

\schritt{Schritt 3 --- Zusammenzählen}
\[ \Pr\bigl[\max \ge a,\ S_n < a\bigr] = \Pr\bigl[\max \ge a,\ S_n > a\bigr] = \Pr[S_n > a] \]
--- denn wer über $a$ endet, hat $a$ zwangsläufig erreicht.

\[ \Pr[\max \ge a] = \Pr[S_n > a] + \Pr[S_n = a] + \Pr[S_n > a] \]
$\blacksquare$

\schritt{Schritt 4 --- Eine Zahl}
Für $n = 2000$ und $a = 50$ liefert die Normalnäherung mit $\sigma = \sqrt{2000} = 44{,}7$:
\[ \Pr[S_n > 50] \approx 1-\Phi(1{,}118) = 0{,}132 \quad\Longrightarrow\quad \Pr[\max \ge 50] \approx 0{,}264 \]

\begin{center}
\begin{minipage}{0.88\textwidth}
\kursiv{\textbf{Ein Viertel aller Verläufe erreicht irgendwann die Fünfzig, obwohl nur ein Achtel dort oben endet.} Das Spiegelungsprinzip ist der Schlüssel zu fast allen Aussagen über Höchststände --- und es wird in Kapitel~10 unverändert auf die Brownsche Bewegung übertragen.}
\end{minipage}
\end{center}
\end{aufgabe}

\begin{aufgabe}
\boxtitel{Aufgabe 1.3 --- Der Spielerruin}
Ein Spieler beginnt mit $a$ Euro, sein Gegner mit $b$. Bei jedem Wurf wechselt ein Euro. Mit welcher Wahrscheinlichkeit verliert er alles?

\loesung
\schritt{Schritt 1 --- Der Ansatz}
Sei $r_k$ die Ruinwahrscheinlichkeit bei Kapital $k$. Dann gilt für $0 < k < a+b$:
\[ r_k = \tfrac12 r_{k-1} + \tfrac12 r_{k+1}, \qquad r_0 = 1,\quad r_{a+b} = 0 . \]

\schritt{Schritt 2 --- Lösen}
Die Gleichung besagt, dass $r_k$ linear in $k$ ist:
\[ r_k = 1 - \frac{k}{a+b} \quad\Longrightarrow\quad r_a = \frac{b}{a+b} \]

\schritt{Schritt 3 --- Die Zahlen}
\begin{center}
{\small
\renewcommand{\arraystretch}{1.46}
\begin{tabular}{@{}P{1.6cm}P{1.6cm}P{2.2cm}P{2.4cm}@{}}
\toprule
$a$ & $b$ & \textbf{Ruin} & \textbf{Spieldauer}\\
\midrule
$10$ & $10$ & $0{,}500$ & $100$\\
$10$ & $100$ & $0{,}909$ & $1000$\\
$10$ & $1000$ & $0{,}990$ & $10\,000$\\
\bottomrule
\end{tabular}}
\end{center}
\kursiv{Die erwartete Spieldauer ist $a\cdot b$ --- man rechnet sie mit demselben Ansatz nach.}

\schritt{Schritt 4 --- Die Lehre}
\begin{center}
\begin{minipage}{0.9\textwidth}
\textbf{Gegen einen unendlich reichen Gegner verliert man bei fairem Spiel mit Sicherheit --- es dauert nur lange.}
\end{minipage}
\end{center}

\begin{center}
\begin{minipage}{0.88\textwidth}
\kursiv{\textbf{Das ist der Grund, warum die Spielbank nicht gewinnen muss, um zu gewinnen} --- ihr Kapitalvorsprung genügt. Und es ist derselbe Mechanismus, der in Kapitel~6 über das Aussterben kleiner Populationen entscheidet: Nicht der Erwartungswert bringt sie um, sondern die absorbierende Schranke bei null.}
\end{minipage}
\end{center}
\end{aufgabe}

\begin{geschichte}
\boxtitelg{Ein Vortrag und ein Irrtum}
\Person{Karl Pearson} fragte 1905 in einem Brief an die Zeitschrift \kursiv{Nature}, wie sich ein Mann verteile, der planlos umherirrt --- und prägte dabei den Namen \enquote{random walk}.

\begin{center}
\begin{minipage}{0.9\textwidth}
\kursiv{\Person{Lord Rayleigh} antwortete eine Woche später, das Problem sei aus der Akustik seit 1880 bekannt. Pearson erwiderte, die Lehre sei: \textbf{\enquote{Der wahrscheinlichste Ort für einen verlorenen Betrunkenen ist irgendwo in der Nähe seines Ausgangspunktes.}} Die Frage war medizinisch gemeint --- es ging um die Ausbreitung von Moskitos in Malariagebieten.}
\end{minipage}
\end{center}

\schritt{Pólya und der Zufall im Park}
\begin{center}
\begin{minipage}{0.9\textwidth}
\kursiv{\Person{George Pólya} kam nach eigener Darstellung auf Satz 1.4, weil er bei Spaziergängen in einem Zürcher Wald einem Paar immer wieder begegnete und sich fragte, wie wahrscheinlich das sei. \textbf{Er veröffentlichte den Satz 1921}; die Unterscheidung zwischen zwei und drei Dimensionen war die Überraschung.}
\end{minipage}
\end{center}

\schritt{Feller und das Arkussinus-Gesetz}
\begin{center}
\begin{minipage}{0.9\textwidth}
\kursiv{\Person{William Feller} stellte Satz 1.3 in seinem Lehrbuch von 1950 an den Anfang --- ausdrücklich, um die Anschauung zu erschüttern. \textbf{Er schrieb, kein anderes Ergebnis der Wahrscheinlichkeitsrechnung widerspreche der Intuition so vollständig}, und empfahl, die Zahlen zu simulieren, bevor man sie glaubt. Genau das tut Abschnitt 1.2 dieses Kapitels.}
\end{minipage}
\end{center}
\end{geschichte}

\begin{forschung}
\boxtitelp{Irrfahrten heute}
\begin{center}
{\small
\renewcommand{\arraystretch}{1.45}
\begin{tabular}{@{}P{3.0cm}P{5.2cm}@{}}
\toprule
\textbf{Auf Graphen} & Mischzeiten, Rückkehrzeiten, Zusammenhang mit Eigenwerten --- Band 16.\\
\textbf{In Zufallsumgebung} & Irrfahrten, deren Übergangswahrscheinlichkeiten selbst zufällig sind.\\
\textbf{Selbstvermeidend} & Pfade ohne Wiederkehr --- Modell für Polymere, bis heute weitgehend ungelöst.\\
\bottomrule
\end{tabular}}
\end{center}

\schritt{Zur dritten Zeile}
\begin{center}
\begin{minipage}{0.88\textwidth}
\kursiv{Die selbstvermeidende Irrfahrt darf keinen Punkt zweimal besuchen. \textbf{Schon die Frage, wie weit sie sich nach $n$ Schritten typischerweise entfernt, ist in drei Dimensionen offen} --- vermutet wird $n^{0{,}588}$, bewiesen ist fast nichts. In zwei Dimensionen kennt man den Exponenten $3/4$ seit 2000 aus der Theorie der konformen Felder; der Beweis brachte \Person{Wendelin Werner} 2006 die Fields-Medaille.}
\end{minipage}
\end{center}
\end{forschung}

\begin{merksatz}
\boxtitel{Was Sie jetzt können}
\begin{center}
\renewcommand{\arraystretch}{1.5}
\begin{tabular}{@{}P{3.0cm}P{5.6cm}@{}}
\toprule
\textbf{Definition 1.1} & Prozess als Familie, Pfad als Realisierung\\
\textbf{Die Wurzel} & typische Entfernung $\sqrt{2n/\pi}$\\
\textbf{Satz 1.3} & Arkussinus --- Ausgeglichenheit ist selten\\
\textbf{Satz 1.4} & Rückkehr sicher bis $d=2$, sonst $0{,}34$\\
\textbf{Spiegelung} & Höchststände aus Endständen\\
\bottomrule
\end{tabular}
\end{center}

\begin{center}
\begin{minipage}{0.88\textwidth}
\kursiv{Die Irrfahrt hat kein Gedächtnis über den gegenwärtigen Ort hinaus. Das nächste Kapitel macht daraus eine Definition --- und findet, dass fast alle brauchbaren Prozessmodelle diese eine Eigenschaft teilen.}
\end{minipage}
\end{center}
\end{merksatz}
